\documentclass[11pt,a4paper]{article}
\usepackage{fontspec}
\setmainfont{DejaVu Sans}
\setmonofont{DejaVu Sans Mono}
\usepackage[margin=2.2cm]{geometry}
\usepackage[french]{babel}
\frenchsetup{AutoSpacePunctuation=false}
\usepackage{amsmath,amssymb}
\usepackage{longtable,array,booktabs}
\usepackage[table]{xcolor}
\usepackage{enumitem}
\usepackage{fancyvrb}
\usepackage{pifont}
\usepackage{hyperref}
\usepackage[normalem]{ulem}
\setlist[itemize]{leftmargin=1.4em,itemsep=2pt,topsep=4pt}
\setlist[enumerate]{leftmargin=1.7em,itemsep=2pt,topsep=4pt}
\newcommand{\code}[1]{\texttt{\detokenize{#1}}}
\newcommand{\fait}{\textcolor{green!45!black}{\ding{51}}}      % ✅
\newcommand{\encours}{\textcolor{blue!65!black}{\ding{46}}}   % 📝
\newcommand{\afaire}{\textcolor{orange!85!black}{$\bigcirc$}}  % pas commencé
\hypersetup{colorlinks=true,linkcolor=black,urlcolor=blue}
\setlength{\tabcolsep}{3pt}
\emergencystretch=2em
\tolerance=2000
\setlength{\LTpre}{6pt}\setlength{\LTpost}{6pt}
\renewcommand{\arraystretch}{1.15}
\pagestyle{plain}
\hypersetup{pdftitle={CDS}}
\begin{document}
\begin{center}{\LARGE\bfseries Cahier des spécifications techniques — v2.1\par}\end{center}\vspace{4pt}
\begin{flushleft}\textbf{Projet} : Reconnaissance d'objets connectés pour audit de sécurité IoT\end{flushleft}

\begin{flushleft}\textbf{Auteur} : Mohamed Salhi — Polytech Lille / ENSA Kénitra · \textbf{Tuteur} : Thomas Vantroys\end{flushleft}

\begin{flushleft}\textbf{Version} : 24/09/2026 · \textbf{Statut} : version de référence — remplace toutes les versions antérieures.\end{flushleft}

\textbf{Compagnon} : le cahier des charges (CDC-v3.1) porte le \textbf{quoi} (périmètre, chiffres mesurés, critères d'acceptation, jalons) ; le présent document porte le \textbf{comment} : comment c'est construit, sur quelles machines, avec quelles garanties.

\begin{quote}\small \textbf{Comment lire ce document.} Écrit \textbf{sans prérequis technique} ; chaque terme utile est expliqué au passage (un lexique complet est dans le CDC §15). Les statuts sont : \fait fait et prouvé · \encours commencé mais incomplet · \afaire pas commencé. Tout chiffre vient d'une mesure datée.\end{quote}
\medskip\hrule\medskip
\section*{1. Les règles de fabrication du logiciel}
\begin{enumerate}  \item \textbf{Les étapes se parlent avec des formulaires fixes.} Chaque étape du traitement transmet à la suivante un objet de type défini (pas un dictionnaire libre où n'importe quoi peut se glisser). Conséquence : un test peut vérifier une étape \textbf{à la seule}, sans rien lancer d'autre.
  \item \textbf{Une base de données simple et locale.} Tout est stocké dans des fichiers SQLite (un format de base de données qui tient dans un fichier, sans serveur à administrer), en mode d'écriture sûre : un seul programme écrit, les autres lisent sans se gêner.
  \item \textbf{Ça tourne sur une seule machine.} La chaîne complète fonctionne sur une machine de 12 Go de mémoire graphique (et même sur le processeur seul pour la partie recherche de fiches/failles). Le cluster de l'école sert à accélérer les tests, il n'est \textbf{jamais indispensable}.
  \item \textbf{Fermeture d'internet pendant l'audit.} Pendant un traitement, \textbf{aucune requête ne sort} vers l'extérieur — prouvé le 16/09 : les sites internet de référence étaient bloqués et l'analyse a quand même tourné.
  \item \textbf{Le code fait foi.} Ce document décrit les choix ; pour le détail d'un algorithme, c'est le code qui dit la vérité (écrit pour être lu).
\end{enumerate}

\medskip\hrule\medskip
\section*{2. Les machines et le modèle d'IA retenus (faits mesurés)}
\subsection*{2.1 Le modèle d'IA}
\begin{itemize}  \item \textbf{Modèle retenu} : \texttt{Qwen2.5-VL-7B-Instruct} en version \textbf{compressée 4 bits} (une compression nécessaire pour tenir dans les 12 Go de la carte graphique — voir le choix détaillé au CDC §3.1).
\end{itemize}

\begin{itemize}  \item \textbf{Comment il est servi} : par le logiciel \textbf{vLLM 0.29.0} qui expose le modèle sur le réseau local (port 8000). Poids : ≈ 6,9 Go ; mémoire réellement utilisée : \textbf{8 962 Mo sur 12 288 Mo} mesurée le 14/09.
\end{itemize}

\begin{itemize}  \item \textbf{Réponses en format rigide} : le modèle est contraint de répondre en \textbf{JSON} (le format d'échange standard des logiciels : paires « clé : valeur », sans phrase autour), température \textbf{0} (aucun hasard), \textbf{640 mots maximum} par réponse — trois réglages \textbf{figés avant les tests}. Résultat : 0 réponse mal formée, 0 réponse coupée en plein milieu sur les 144 photos.
\end{itemize}

\begin{itemize}  \item \textbf{Rappel du choix} (détail et corrections datées au CDC §3.1) : 79/144 pour le modèle retenu contre 53 pour la génération suivante (écart réel), et un modèle concurrent écarté car il respecte mal le format (21 réponses sur 144 hors format).
\end{itemize}

\subsection*{2.2 Les machines de l'école}
\begin{itemize}  \item Trois machines identiques : processeur i9, 62 Go de RAM, carte graphique \textbf{RTX 3060 de 12 Go}, système sans gestionnaire de services (tout se lance à la main, par port), pilotes graphiques \textbf{alignés entre les trois}, recette de contrôle \textbf{10/10} (15/09).
\end{itemize}

\begin{itemize}  \item \textbf{État au 24/09} : relevés du jour — 20h56 une machine injoignable → 21h11 deux machines en service → \textbf{21h59 les trois en service}, temps dégradé 2,1 s / 2,1 s / 5,1 s sur 4 photos, \textbf{réponses identiques entre les trois machines}.
\end{itemize}

\begin{itemize}  \item \textbf{Ces machines sont partagées avec d'autres étudiants} : à un moment, l'une d'elles était saturée par un logiciel tiers (11,6 Go sur 12 Go de mémoire graphique) et un réseau privé tiers était actif dessus. \textbf{Décision consignée (24/09)} : périmètre contractuel = une machine saine ; les trois ne sont qu'un banc d'accélération.
\end{itemize}

\medskip\hrule\medskip
\section*{3. L'organisation du traitement, étape par étape}
Chaque flèche ci-dessous est un \textbf{type de données fixe}, pas une boîte noire :

\begin{Verbatim}[fontsize=\small]
[Photo] ──► 1. Lecture par le modèle de vision (+ lecture du texte)
                 │  ▼ sortie brute du modèle (marque, modèle, type, texte lu)
             2. Consolidation (accord entre photos × vérification officielle)
                 │  ▼ identification consolidée + statut
             3. Résolution de fiche produit (recherche approximative)
                 │  ▼ candidat(s) CPE + score
             4. Recherche de failles connues (locales, en Python)
                 │  ▼ liste de CVE (+ mention « niveau produit »)
             5. Rapport : journal scellé + PDF
\end{Verbatim}
\subsection*{3.1 Les jeux de valeurs fermés}
\sloppy
\begin{longtable}{|p{0.200\linewidth}|p{0.711\linewidth}|}
\hline
\textbf{Domaine} & \textbf{Valeurs autorisées} \\ \hline
\endhead
\textbf{Type d'objet} & prise, caméra, ampoule, routeur, enceinte \textbf{+ \texttt{other}} (le « je ne sais pas ») \\ \hline
\textbf{Statut d'identification} & identifié / identifié avec réserve / non identifié — \textbf{trois valeurs, jamais deux} (un système qui ne peut que répondre « oui » ne peut pas reconnaître son ignorance) \\ \hline
\textbf{Niveau de preuve} & correspondance exacte au produit / plage de version / simple vendeur — exposé tel quel à l'auditeur (§9.3) \\ \hline
\textbf{D'où vient l'étiquetage} & catégorie officielle / nom de fichier / annonce de vente / saisie manuelle — traçable pour détecter les fuites de réponse \\ \hline
\end{longtable}
\subsection*{3.2 Les formulaires échangés entre étapes}
\sloppy
\begin{longtable}{|p{0.180\linewidth}|p{0.731\linewidth}|}
\hline
\textbf{Élément transmis} & \textbf{Contenu essentiel} \\ \hline
\endhead
\textbf{Sortie du modèle de vision} & marque, modèle, type d'objet, texte lu sur l'objet, « confiance ressentie » par le modèle (\textbf{jamais affichée} à l'utilisateur : elle n'est pas fiable) \\ \hline
\textbf{Candidat de fiche produit} & code CPE, fabricant, produit, version, scores, méthode employée \\ \hline
\textbf{Résultat de faille} & identifiant de la faille, sévérité, date, \textbf{niveau de preuve} (produit exact / plage de version / vendeur entier) \\ \hline
\textbf{Identification finale} & marque, modèle, statut, nombre de photos, \textbf{accord} et \textbf{résolvabilité} (§8), distribution des réponses, fiches et failles associées \\ \hline
\textbf{Provenance du traitement} & empreintes numériques (empreinte = empreinte digitale d'un fichier) des poids du modèle, de la consigne et du schéma de réponse ; versions du logiciel ; \textbf{empreinte de la base des failles} ; dates de début et de fin \\ \hline
\textbf{Événement de journal} & numéro de ligne \textbf{toujours croissant} (résiste à une horloge qui dérive), étape, message, charge utile, horodatage \\ \hline
\textbf{Entrée photo} & identifiant, empreinte, provenance, licence, date de collecte, dimensions, présence ou non de texte lisible \\ \hline
\end{longtable}
\subsection*{3.3 Ce que le modèle doit obligatoirement renvoyer}
Un objet JSON à champs fixes : \texttt{brand}, \texttt{model}, \texttt{category}, \texttt{visible\_\allowbreak{}text}, \texttt{identification\_\allowbreak{}basis} — \textbf{aucun champ supplémentaire toléré}, longueur maximale par champ, type d'objet pris dans la liste fermée (§3.1). Ce format garantit une réponse \textbf{lisible par la machine}, pas son exactitude.

\medskip\hrule\medskip
\section*{4. Les quatre bases de données locales}
\sloppy
\begin{longtable}{|p{0.148\linewidth}|p{0.590\linewidth}|p{0.159\linewidth}|}
\hline
\textbf{Base} & \textbf{Ce qu'elle contient réellement} & \textbf{État} \\ \hline
\endhead
\textbf{Photos du projet} (500 Ko) & 48 objets, 144 photos utilisables + 11 retirées, 288 relevés de lecture de texte, 2 traitements enregistrés ; les tables de résultats bruts et de cas de calibration existent mais sont \textbf{encore vides} & \fait vérité terrain gelée \\ \hline
\textbf{Dictionnaire des produits CPE} (463 Mo) & \textbf{1 838 235 fiches} au total, dont \textbf{64 842} concernent du matériel (le périmètre du projet) ; empreinte d'identification enregistrée (14/09) & \fait \\ \hline
\textbf{Base des failles NVD} (154 Mo, importée le 24/09) & \textbf{121 719 failles} et \textbf{324 470 critères matériels} répartis sur \textbf{873 fabricants} ; les 4 fichiers d'origine sont tracés par empreinte ; statuts réellement présents : 83 617 modifiées, 23 569 analysées, 10 855 différées, \textbf{3 559 retirées par le NVD} (et conservées telles quelles, §7.6), 119 diverses & \fait premier import — les preuves complètes du §7 restent \afaire \\ \hline
\textbf{Référentiel de vérification du texte} & un jeu public annoté (COCO-Text) servant à noter notre outil de lecture & \fait \\ \hline
\end{longtable}
Toutes les photos publiables portent leur \textbf{adresse d'origine, leur licence (CC-BY / CC-BY-SA) et leur empreinte} ; les fichiers d'images ne sont \textbf{pas diffusés} avec le rapport.

\medskip\hrule\medskip
\section*{5. La lecture automatique du texte (OCR) — un aide, pas une vérité}
\begin{itemize}  \item \textbf{Deux logiciels} comparés côte à côte (PaddleOCR et OnnxTR) ; on retient un texte si \textbf{l'un ou l'autre} le lit (choix du 15/09 : 50 photos avec texte, 94 sans, sur 144).
\end{itemize}

\begin{itemize}  \item \textbf{Trois mesures, trois situations} — lisons-les du plus favorable au plus réaliste :
\end{itemize}

\sloppy
\begin{longtable}{|p{0.209\linewidth}|p{0.295\linewidth}|p{0.392\linewidth}|}
\hline
\textbf{Situation} & \textbf{Résultat} & \textbf{En clair} \\ \hline
\endhead
On dessine le rectangle autour du texte & \textbf{66 \%} de réussite (1 384 textes, fourchette 63–68 \%) & il recopie bien les 2 textes sur 3 quand on lui montre où chercher \\ \hline
Photo entière, gros plan & \textbf{35 \%} (600 photos, 2 604 mots, fourchette 32–38 \%) & il lit environ 1 mot sur 3 \\ \hline
Vraie vie : trouver puis lire & \textbf{10 \%} (12 sur 120, fourchette 6–17 \%) & \textbf{la chaîne complète est faible} — on n'a pas le droit de compter dessus \\ \hline
\end{longtable}
\begin{itemize}  \item \textbf{Rôle} : alimenter le champ « texte lu » et le drapeau « photo contient du texte » ; remplacer \textbf{jamais} l'étiquetage humain. Quand un modèle affirme avoir lu une étiquette, le système vérifie que le texte est bien présent (critère AC2 du CDC : 94 \% de conformité mesurée).
\end{itemize}

\medskip\hrule\medskip
\section*{6. Retrouver la bonne fiche produit (résolution CPE)}
\begin{itemize}  \item \textbf{Comment ça marche} : on cherche d'abord le \textbf{fabricant exact} (petit dictionnaire d'alias pour les orthographes : « tp-link », « tplink », « tp link »), puis on classe les produits par \textbf{similarité de mots} (fonction de rapprochement qui ignore les mots superflus), puis on décide par \textbf{marge} : si le meilleur candidat n'est pas nettement devant le deuxième, on refuse de trancher.
\end{itemize}

\begin{itemize}  \item \textbf{Ses seuils} : acceptation automatique à \textbf{0,90} ; plancher à \textbf{0,60} (en dessous : refus explicite, jamais de réponse approximative) ; marge minimale \textbf{0,08} entre le 1\textsuperscript{er} et le 2\textsuperscript{e}. Méthode de réglage : on balaie les seuils et on garde \textbf{le plus souple qui garde une fiabilité ≥ 95 \%}. \textbf{Ce réglage définitif attend le banc de 200 cas (porte G2).}
\end{itemize}

\begin{itemize}  \item \textbf{Index de recherche rapide (FTS5)} : prévu pour retrouver un produit même mal orthographié — \textbf{pas encore en place} (constat du 24/09) \afaire.
\end{itemize}

\begin{itemize}  \item \textbf{État actuel du banc de mesure} : sur 68 cas volontairement dégradés : \textbf{37 \% au 1\textsuperscript{er} choix, 56 \% dans le top 3}, contre des cibles de 80 \% et 92 \% — \textbf{échec public} (critère AC10). Causes diagnostiquées : des produits grand public \textbf{absents} du dictionnaire, deux orthographes concurrentes sans table d'alias, et un refus par marge volontairement prudent. \textbf{Levier identifié} : un tableau d'alias manuel d'environ 50 entrées.
\end{itemize}

\begin{itemize}  \item \textbf{Refus explicites} : \texttt{E\_\allowbreak{}NO\_\allowbreak{}CPE\_\allowbreak{}MATCH} (aucune fiche trouvée) et \texttt{E\_\allowbreak{}CPE\_\allowbreak{}AMBIGUOUS} (plusieurs fiches plausibles) — le système présente ses doutes plutôt qu'une fausse certitude.
\end{itemize}

\medskip\hrule\medskip
\section*{7. La base locale des failles (NVD) — première importation faite le 24/09}
\subsection*{7.1 D'où viennent les données et comment on les trace}
Les fichiers officiels annuels du NVD sont téléchargés une fois, conservés sur disque, et \textbf{signés par empreinte} (sha256 — l'empreinte digitale du fichier) : nom, empreinte, date de publication, date d'import, nombre de lignes et de failles. C'est cette traçabilité qui rend un rapport reproductible : le rapport dit exactement \textbf{quel fichier de failles} il a consulté.

\subsection*{7.2 Ce qui est conservé pour chaque faille}
L'identifiant (ex. CVE-2024-1234), son \textbf{statut officiel brut} (dont les \textbf{3 559 failles retirées} par le NVD — on les garde pour ne jamais se demander « pourquoi celle-là a disparu ? »), les dates, la description, et la sévérité avec son score de criticité.

\subsection*{7.3 Les critères d'association à un produit}
Une ligne par association : le \textbf{motif officiel original} conservé tel quel (c'est la pièce montrant \textit{pourquoi} cette faille remonte), le fabricant, le produit, la version, les quatre bornes de plage (début/fin, inclus/exclus), les drapeaux « vulnerable », « exclure » et l'opérateur logique — avec interdiction de doublons.

\subsection*{7.4 On ne garde que le matériel}
Seules les associations vers du \textbf{matériel} (\texttt{part='h'}) sont importées — \textbf{64 842 entrées matérielles sur 1 838 235} au total. Le reste (logiciels, systèmes d'exploitation) est écarté au fil du téléchargement, sans jamais charger une année entière en mémoire.

\begin{quote}\small \textbf{Ce que ce filtre coûte, mesuré le 29/09/2026} : pour les vendeurs du corpus du projet, la base ne garde que \textbf{10 failles sur 11 714} (0,09 \%) — TP-Link, D-Link, Hikvision, AVM, Philips, Zyxel, MikroTik : \textbf{zéro} faille matérielle sur 2021–2024, alors que les mêmes produits en ont des centaines classées \texttt{part='o'} (système) ou \texttt{part='a'} (application). Voir CDC §14.3.10 et \texttt{benchmarks/cve.md} §2.\end{quote}
\subsection*{7.5 On peut relancer l'import sans risque}
Si le schéma existe déjà, on ne crée rien deux fois ; l'écriture ignore les doublons. \textbf{Relancer l'import doit donc produire exactement les mêmes comptes} — c'est la preuve exigée par la porte G1, et elle est \textbf{faite le 29/09}. Un compteur de progression permet de reprendre au milieu d'un fichier si le téléchargement est interrompu.

\textbf{Nuance apportée le 29/09 :} relancer le \textbf{même} fichier ne prouve que la \textit{provenance} — son empreinte le fait sauter avant toute écriture, donc le chemin d'écriture n'est jamais exécuté. Le banc rejoue donc les flux une \textbf{troisième} fois sous un autre nom, ce qui force l'écriture : les comptes de données restent identiques et il n'y a \textbf{aucun doublon}. Les deux passes sont décrites dans \texttt{benchmarks/cve.md} §3.

\subsection*{7.6 Les failles retirées restent visibles}
Si le NVD retire une faille, elle est \textbf{conservée avec son statut} : un rapport antérieur qui la mentionne garde son explication. La recherche les met à part et les signale, sans les mélanger aux failles actives.

\begin{quote}\small \textbf{État (24/09, 21h36)} : \textbf{premier import exécuté} — 4 années (2021–2024), \textbf{121 719 failles / 324 470 critères matériels en 11 secondes}, empreinte par fichier. \textasciitilde{}\textasciitilde{}Porte G1 non fermée : deuxième import sans doublon et vérification à la main de 60 cas \afaire ; le temps de recherche n'a été mesuré que sur 6 requêtes (3,5 ms).\textasciitilde{}\textasciitilde{} \textit{(périmé — voir l'état du 29/09 ci-dessous)}\end{quote}
\begin{quote}\small \textbf{État (29/09, 13h50)} — banc \texttt{benchmarks/run\_\allowbreak{}cve.py}, résultat complet dans \texttt{benchmarks/cve.json} : - \textbf{rejeu sans doublon} : 3 passes sur une base neuve (flux d'origine, mêmes noms, copies renommées) ⇒ comptes de données \textbf{identiques}, \texttt{COUNT(*) = COUNT(DISTINCT …)}, \textbf{0 doublon} ✓ ; - \textbf{60 cas} ⇒ \textbf{0 faux négatif, 0 faux positif} ✓ — vérité terrain re-dérivée des \texttt{.json.gz} \textbf{bruts} par un lecteur indépendant de l'ingestion ; - \textbf{latence} : p95 \textbf{0,3–0,6 ms} sur 60 requêtes distinctes (600 exécutions), cible 50 ms ✓ ; - \afaire \textbf{reste la relecture humaine des 60 cas} — et c'est précisément là que la décision \texttt{part='h'} doit être tranchée (CDC §14.3.10). - correction jointe : \texttt{feed\_\allowbreak{}versions.row\_\allowbreak{}count} annonçait le nombre de critères \textbf{produits} au lieu du nombre de lignes \textbf{insérées} ; corrigé dans le code et dans les quatre lignes déjà écrites (CDC §14.3.8).\end{quote}
\medskip\hrule\medskip
\section*{8. La confiance : deux vérifications croisées (et non une seule)}
\begin{itemize}  \item \textbf{1\textsuperscript{er} vérification — l'accord} : le modèle est interrogé sur \textbf{plusieurs photos du même objet} (prises de côtés différents). S'il répond la même chose à chaque fois, il est « d'accord avec lui-même ». C'est utile, mais insuffisant : un modèle qui a mémorisé une réponse fausse sera d'accord avec lui-même aussi.
\end{itemize}

\begin{itemize}  \item \textbf{2\textsuperscript{e} vérification — la résolvabilité} : la réponse consensuelle \textbf{existe-t-elle réellement} dans le dictionnaire officiel des produits ? Si la marque renvoyée n'existe dans aucune fiche, la réponse est inventée.
\end{itemize}

\begin{itemize}  \item \textbf{Le cas dangereux} : accord très élevé \textbf{et} résolvabilité faible = \textbf{réponse inventée, mais répétée avec assurance}. C'est précisément le cas le plus fréquent observé sur nos tests — d'où la règle : ce cas est \textbf{refusé} (statut « non identifié »), jamais affiché comme une réussite.
\end{itemize}

\begin{itemize}  \item \textbf{La confiance que le modèle affirme lui-même n'est jamais montrée à l'utilisateur} : elle n'est pas fiable (écart mesuré entre ce qu'il déclare et ce qu'il fait : déclaré − juste = \textbf{−32 points} sur le banc de calibration).
\end{itemize}

\begin{itemize}  \item \textbf{Aucun score de risque agrégé} ; chaque faille remontée porte la mention « niveau produit » (§9.3).
\end{itemize}

\medskip\hrule\medskip
\section*{9. Chercher les failles connues d'un produit}
\textbf{Validé le 24/09} : \textbf{40 tests automatiques passent} ; sur la vraie base de 324 470 critères, une recherche prend \textbf{1,5 ms en médiane, 3,5 ms en dégradé} (6 essais) — largement sous la cible contractuelle de 50 ms.

\textbf{Mesuré le 29/09} sur \textbf{60 requêtes distinctes} (et 600 exécutions) : p50 \textbf{0,05–0,13 ms}, p95 \textbf{0,3–0,6 ms}, max 5,2 ms — cible de 50 ms tenue avec trois ordres de grandeur de marge. Les 60 cas et le rejeu sans doublon sont produits par \texttt{benchmarks/run\_\allowbreak{}cve.py} ; il reste la relecture \textbf{humaine} des 60 cas, et la décision de périmètre \texttt{part='h'} (CDC §14.3.10) qu'aucun test ne peut trancher.

\subsection*{9.1 Le tri préalable (en base de données) ne tranche rien tout seul}
Une première requête élimine l'essentiel et ramène des \textbf{candidats larges} (fabricant exact, produit exact \textbf{ou} tout produit du fabricant) ; la décision finale se prend ensuite en Python, sur des fonctions simples qu'on peut tester une par une — plutôt que dans une requête compliquée que personne ne pourrait relire.

\subsection*{9.2 Les plages de versions, évaluées en Python}
Une faille concerne souvent une plage (« de 1.0 à 2.0 exclu »). La comparaison des versions se fait \textbf{en Python} et de façon lisible : comparaison \textbf{numérique} (donc « 1.10 » est bien plus récent que « 1.9 ») et une version d'essai (« 1.2.3-beta ») est considérée comme \textbf{antérieure} à sa version finale (« 1.2.3 »).

\textbf{Le cas nominal du projet} : nous ne connaissons \textbf{pas} la version du logiciel interne (une photo ne le donne pas). Le résultat est alors marqué \textbf{\texttt{unresolved}} : « non résolu ». Il n'est ni déclaré concerné, ni déclaré non concerné \textbf{en silence} — c'est la seule manière de ne pas fabriquer de faux négatifs, le défaut le plus grave d'un tel outil.

\subsection*{9.3 Niveau de preuve, affiché tel quel}
Trois niveaux, jamais confondus : \textbf{produit exact} (preuve forte) > \textbf{plage de version} (preuve forte) > \textbf{vendeur entier} (preuve faible, quand le critère ne nomme que le fabricant). L'auditeur voit ce niveau dans le résultat ; un match faible n'est jamais fusionné avec un match fort. Chaque faille porte \textbf{\texttt{triage\_\allowbreak{}scope='product'}} — le champ qui, dans la base comme dans le PDF, interdit d'interpréter le résultat au niveau de l'exemplaire.

\medskip\hrule\medskip
\section*{10. Le journal scellé et la signature (intégrité du rapport)}
\begin{itemize}  \item \textbf{Déjà prévu dans les formulaires} : la provenance complète (§3.2) et la ligne de journal avec numéro \textbf{toujours croissant} — le journal survit à une horloge qui dérive.
\end{itemize}

\begin{itemize}
  \item \textbf{À construire (\afaire, portes G3/G4)} :
\begin{enumerate}  \item \textbf{chaîne de scellement} : chaque ligne du journal est scellée avec la précédente (HMAC, une fonction cryptographique) — modifier une ligne casse toute la suite, la falsification est détectée à la ligne exacte ;
  \item \textbf{signature de la dernière ligne} par une clé privée ; la clé publique est \textbf{imprimée dans le rapport} : un tiers vérifie le document \textbf{sans recevoir aucun secret} ;
  \item l'empreinte de la base des failles est écrite dans chaque traitement ;
  \item \textbf{6 tests d'altération} dont deux indispensables : un test de \textbf{fausse ligne créée avec la vraie clé, qui doit réussir} (la limite est démontrée, pas cachée) et un test « rapport modifié après coup ⇒ refus ».
\end{enumerate}
\end{itemize}

\begin{itemize}  \item \textbf{Formulation exacte à défendre} : le scellement prouve que le contenu \textbf{n'a pas changé après} la signature ; il ne prouve pas à lui seul l'identité du signataire — c'est le rôle de la signature à clé publique de la dernière ligne.
\end{itemize}

\medskip\hrule\medskip
\section*{11. Le rapport PDF et l'interface (\afaire)}
\begin{itemize}  \item \textbf{Contenu du PDF} : la photo, la fiche produit retenue avec ses candidats et scores, les failles du modèle avec leur niveau de preuve, la provenance complète, et \textbf{quatre mentions obligatoires} : triage au niveau produit, logiciel interne non déterminé, message non rassurant si aucune faille, date et empreinte de la base des failles utilisée.
\end{itemize}

\begin{itemize}  \item \textbf{Interface de démonstration} : une page web simple \textbf{ou} une commande en ligne — parcours complet exigé par la porte G5 (07/12) : envoyer une photo → lancer le traitement → récupérer le PDF → \textbf{le vérifier hors-ligne} avec la clé publique.
\end{itemize}

\medskip\hrule\medskip
\section*{12. Temps de réponse mesurés et cas dégradés}
\textit{Note de lecture} : le temps \textbf{dégradé} (p95) est le temps en deçà duquel se situent 95 requêtes sur 100 — c'est lui qu'on compare à la cible.

\sloppy
\begin{longtable}{|p{0.321\linewidth}|p{0.308\linewidth}|p{0.268\linewidth}|}
\hline
\textbf{Mesure} & \textbf{Valeur} & \textbf{Situation de mesure} \\ \hline
\endhead
Temps par image, médian / dégradé & 1,08 s / \textbf{1,88 s} & 12 mesures, 15/09 \\ \hline
Temps dégradé en vraie chaîne (photo jamais vue) & \textbf{3,9 s} (max 5,5 s) & flux réel — la valeur honnête \\ \hline
Débit sur les 3 machines ensemble & \textbf{57,9 requêtes/s} & banc de charge du 15/09 \\ \hline
Panne d'un nœud en plein traitement & \textbf{2400 requêtes sur 2400 servies}, 0 abandon & mode dégradé déjà testé \\ \hline
Banc des 3 machines (24/09, 4 photos) & dégradé 2,1 / 2,1 / 5,1 s, réponses identiques & les trois en inférence réelle \\ \hline
\textbf{Cible contractuelle} & dégradé \textbf{< 10 s} & \fait tenue avec une marge ×2,6 en flux réel \\ \hline
\end{longtable}
Les temps \textbf{ne sont jamais comparés entre protocoles différents sans le dire} (CDC §3.1).

\textbf{Modes dégradés prévus} : base des failles indisponible → identification seule, message explicite ; machine perdue en cours de traitement → traitement partiel comptabilisé comme tel ; aucune faille trouvée → message non rassurant (\texttt{E\_\allowbreak{}NO\_\allowbreak{}CVE}).

\medskip\hrule\medskip
\section*{13. Licences, données personnelles, conformité}
\begin{itemize}  \item Photos sous licences \textbf{CC-BY / CC-BY-SA} (Wikimedia Commons), avec adresse d'origine, licence et empreinte pour chacune ; \textbf{jamais diffusées} avec le livrable.
\end{itemize}

\begin{itemize}  \item Aucune donnée personnelle dans le corpus ; \textbf{aucun scan réseau}, aucune collecte de terrain (hors projet — articles 323-1 et suivants du Code pénal).
\end{itemize}

\begin{itemize}  \item \textbf{Documents de conformité à produire avant la soutenance} : durée de conservation des données, principe de minimisation écrit — \afaire.
\end{itemize}

\begin{itemize}  \item Les données de localisation (EXIF) sont \textbf{retirées dès l'import}, avant tout stockage.
\end{itemize}

\medskip\hrule\medskip
\section*{14. Sécurité d'exploitation (en clair)}
\begin{itemize}  \item \textbf{Fermeture totale d'internet pendant un traitement} : règles de pare-feu écrites \textbf{en mémoire uniquement} (rien n'est écrit sur le système), autorisant seulement le réseau local, le serveur de noms et la boucle locale ; un \textbf{garde de 300 secondes} remet tout dans l'état automatiquement ; le script de restauration est exécuté \textbf{avant} d'appliquer quoi que ce soit. Appliqué puis restauré sur les trois machines le 16/09 — incident d'un réseau tiers détecté et réparé le jour même.
\end{itemize}

\begin{itemize}  \item \textbf{Mots de passe} : transmis par un canal indirect à deux sauts (poste → passerelle → machine), \textbf{jamais dans une ligne de commande} (visible par tous avec la commande \texttt{ps}) \textbf{ni dans un fichier}.
\end{itemize}

\begin{itemize}  \item \textbf{Pas de gestionnaire de services} sur ce système : tout se lance et se vérifie à la main, par port et par numéro de processus, avec un script de recette automatique (version rapide + version approfondie).
\end{itemize}

\medskip\hrule\medskip
\section*{15. Mesure, notation et reproductibilité}
\subsection*{15.1 Un barème figé, un seul outil de notation}
Le barème de notation des photos est \textbf{figé avant} toute mesure ; l'outil qui l'applique est \textbf{unique} et le même partout — \textbf{jamais recopié à la main}, jamais réécrit « de mémoire ».

\subsection*{15.2 Par quoi on compte}
Les marges d'erreur (fourchettes à 95 \%) ne s'appliquent qu'au niveau \textbf{objet} (48) : huit photos du même objet ne sont pas huit épreuves indépendantes. \textbf{Aucune affirmation de précision n'est validée sans fourchette au niveau objet} ; les chiffres comptés par photo ne sortent qu'avec la mention « indicatif ». Toute figure indique son unité de comptage.

\subsection*{15.3 L'incident fondateur (conservé pour la mémoire)}
Une \textbf{copie} de l'outil de notation, refaite « de mémoire », avait publié \textbf{77 \% au lieu de 45 \%} — 32 points trop haut — parce qu'elle écrasait les séparateurs lors du nettoyage des textes. Conséquence : la correspondance des outils est aujourd'hui \textbf{vérifiée automatiquement} (l'outil doit reproduire exactement les chiffres déjà publiés), pas relue à l'œil.

\subsection*{15.4 Ce qui est déjà reproductible}
Deux exécutions complètes d'une heure d'intervalle ont donné des réponses \textbf{identiques à la première} (144/144) ; les réglages (640 mots maximum, garde de fin de texte) ont été \textbf{figés avant les mesures}.

\subsection*{15.5 La bande de sensibilité (règle anti-complaisance)}
Chaque mesure est rejouée sous \textbf{4 barèmes de notation différents} (le plus permissif en 5\textsuperscript{e} point de la bande) : la \textbf{direction de la conclusion doit être la même sous les 4}. Les valeurs sont publiées \textbf{côte à côte} (borne la plus sévère et borne la plus souple), jamais fusionnées en un chiffre unique.

\subsection*{15.6 État restant}
\begin{itemize}  \item Figures principales \textbf{reproductibles} \fait (outil unique + manifestes figés).
\end{itemize}

\begin{itemize}  \item \textbf{Un ancien fichier d'agrégats n'est pas régénérable} (aucun script producteur conservé) — ❌ à réparer avant la porte G6 (14/12).
\end{itemize}

\begin{itemize}  \item Tables de résultats bruts du corpus : \textbf{vides} — les résultats publiés vivent dans les archives de mesure, pas encore en base.
\end{itemize}

\medskip\hrule\medskip
\section*{16. Annexe — les formules utilisées}
\textbf{Marges d'erreur (fourchette de Wilson, confiance 95 \%)} — utilisée pour tout pourcentage publié :

$$\hat{p} \pm \tilde{\Delta} \quad\text{avec}\quad \tilde{\Delta} = \frac{z}{1 + \frac{z^2}{n}}\sqrt{\frac{\hat{p}(1-\hat{p})}{n} + \frac{z^2}{4n^2}} \quad (z = 1{,}96)$$

\textit{Lecture} : \texttt{p̂} est le pourcentage observé, \texttt{n} l'effectif. Cette formule est préférée au calcul naïf parce qu'elle \textbf{reste valable pour de petits effectifs et près de 0 \% ou 100 \%} — exactement notre cas.

\textbf{La confiance du projet, en deux comptages :}

\begin{itemize}  \item \textbf{accord} = proportion de photos donnant la réponse la plus fréquente : $A = \dfrac{\max_t |\{i : t_i = t\}|}{N}$ (publié en fraction brute, ex. « 6 sur 8 », sans marge d'erreur) ;
\end{itemize}

\begin{itemize}  \item \textbf{résolvabilité} = meilleur score de rapprochement du consensus dans le dictionnaire officiel ;
\end{itemize}

\begin{itemize}  \item \textbf{décision} : accord élevé \textbf{et} résolvabilité faible ⇒ réponse inventée et répétée ⇒ \textbf{refus explicite}.
\end{itemize}

\end{document}
